\runninghead{A Kingdom from Above}
\properlane{kingdom-not-from-hence}{A Kingdom from Above, Gathered from Every Nation and Served by Kings}

Christ's kingdom takes neither its origin nor its means from this world, and
for that reason it reaches everywhere in the world and claims everyone in it. Its subjects are those who hear the truth, gathered from every
nation. The kings and nations whose subjection the Collect asks and whose
service the Gradual sings are brought under a rule that works by truth and
grace. That rule takes no earthly throne for itself, and yet rulers owe it, as
rulers, the service of their own office. Their subjection begins in free
conversion now and is completed at the judgment.

\subsection{A prisoner who claims the nations}
\label{sec:hence-question}

The Gospel and the chants around it seem to pull apart. In the governor's
hall, bound and handed over by his own nation, Jesus says three times in two
sentences that his kingdom does not belong to the world around him: it is not
\latin{de hoc mundo}, it would be defended by his servants if it were
\latin{ex hoc mundo}, and it is not \latin{hinc}, from here. Yet the Gradual
before the Gospel gives him dominion ``from sea to sea, and from the river
unto the ends of the earth'' and sings that ``all kings of the earth shall
adore him: all nations shall serve him.'' The Offertory after it has the Father
promise him ``the utmost parts of the earth for thy possession,'' in a psalm
whose next verses bid the kings ``understand'' and ``Serve ye the Lord with
fear'' (Ps 2:10--11). The Collect asks that all the families of the nations be
subjected, \latin{subdantur}, to his rule. And the Postcommunion has the
faithful \latin{militare}, fight as soldiers, under his standards, where the
Gospel's King says that his servants do not fight. How can a kingdom that is
not of this world, and not from hence, claim the whole earth, the nations as
inheritance and the service of every king?

Either side, read without the other, invites a misreading. The
Gospel's denial, read alone, could empty the Collect's petition for the
nations' subjection of any public consequence. \work{Quas primas}, which
instituted the feast, names this error itself: it calls it a grave error to
say that Christ has no authority whatever in civil affairs, and draws from the
feast a public duty of rulers towards him (nn.~17--19, 32). The Gradual and the
Collect, read alone, could make Christ a rival of Caesar for his throne, and
Pilate's first question, whether he is the king of the Jews, already raises
that suspicion.

\subsection{St Augustine: not of this world, but in it}
\label{sec:hence-augustine}

St Augustine begins from Jesus' counter-question, ``Sayest thou this thing of
thyself?'' It was put, he says, so that the vain notion men had formed of his
kingdom, Gentiles and Jews alike, might be shown up (\work{Tractates on the
Gospel of John} 115.1). Then he gives the answer its widest audience: ``Hear
then, ye Jews and Gentiles; hear, O circumcision; hear, O uncircumcision; hear,
all ye kingdoms of the earth: I interfere not with your government in this
world'' (115.2). The reassurance rests on two exact distinctions in the
Lord's words. He did not say that his kingdom is not \emph{in} this world, but
that it is not \emph{of} this world; and he did not say that it is not
\emph{here}, but that it is not \emph{from hence}. ``For His kingdom is here
until the end of the world, having tares intermingled therewith until the
harvest,'' and yet it is not from here, ``for it only sojourns as a stranger in
the world.'' The kingdom itself, for Augustine, is its people, ``Those who
believe in Him,'' who were of the world while they belonged to its prince and
have been made a kingdom no longer of the world by being reborn in Christ. Augustine gives the Epistle's verse as
the account of that change: ``For so did God rescue us from the power of
darkness, and translate us into the kingdom of the Son of His love'' (Col
1:13), and he hears the king set on Sion in the second psalm speaking of the
same hill, which ``is not of this world'' (Ps 2:6).

``Thou sayest that I am a king'' is balanced, he goes on, so that Jesus neither
denies that he is a king nor confesses himself such a king as Pilate imagines:
``Carnal thyself, thou sayest it carnally'' (115.3). And the last words of the
pericope define the subjects of this kingdom. ``Every one that is of the
truth heareth my voice'': hearing is obeying and believing, and in this word
Christ ``commendeth the grace whereby He calleth according to His own
purpose'' (115.4).

\subsection{St John Chrysostom and St Thomas: an empire from above}
\label{sec:hence-chrysostom}

St John Chrysostom reads the same sentence as a statement about the kingdom's
power rather than its people. Is his kingdom then not of this world also?
``Certainly it is. How then says He it `is not'? Not because He does not rule
here, but because He has his empire from above, and because it is not human,
but far greater than this and more splendid.'' The proof is the servants who do
not fight: ``Here He shows the weakness of kingship among us, that its strength
lies in servants; but that which is above is sufficient for itself, needing
nothing.'' He says his kingdom is not from hence ``not depriving the world of
His providence and superintendence, but showing \dots\ that His power was not
human or perishable'' (\work{Homilies on John} 83.4). His paraphrase of the
whole reply is ``I am indeed a King, yet not such an one as you suspect, but far
more glorious'' (84.1). St Cyril of Alexandria makes the same reading a
reassurance to Pilate: Christ did not deny that he was King, but showed that he
was no enemy to Caesar, for he reigns as God over heaven and earth, ``as God He
is King, whether man so will, or no'' (\work{Commentary on John} XII).

St Thomas Aquinas sets the two readings side by side. \latin{Regnum} in the
Gospel may mean the people who are ruled, as Augustine takes it, or the royal
power itself, as Chrysostom does, whose origin is not from worldly causes; and
he gives the Alleluia's verse for that power, an everlasting power that shall
not be taken away. Then he concludes: \latin{non est hinc \dots\ est tamen hic, quia ubique
est}, it is not from here, and yet it is here, because it is everywhere; and he
cites for it Wisdom's reaching ``from end to end'' (Wis 8:1), the Offertory's
verse, and the first half of Daniel's (\work{Super Io.}\ 18 lect.~6). The
point can be put in terms of origin: a kingdom whose origin is one place in
the world is bounded by that place, while one whose origin is above is bounded
by no place, and so it can be present everywhere.

\subsection{The kings of the Gradual and of the second psalm}
\label{sec:hence-kings}

Aquinas's answer says where the kingdom is. What the kings and nations of the
Gradual and the Offertory are to do, the commentators on the two psalms say
concretely: they speak of kings who lead men into the Church, of laws, of free
conversion, and of a homage that some give now and all give at the
judgment.

St Augustine reads the Gradual's psalm of the Church spread through the whole
world, and when he comes to the gifts the kings of Tharsis and Saba ``shall
lead,'' he asks why the verb is \emph{lead}, a word for things that walk. The
gifts ``seem to me to signify men, whom into the fellowship of the Church of
Christ the authority of kings doth lead: although even persecuting kings have
led gifts, knowing not what they did, in sacrificing the holy Martyrs''; and
at once he quotes the Gradual's verse, ``And there shall adore Him all kings of
the earth, all nations shall serve Him'' (\work{Enarrationes in Psalmos}
71.13). The reason for that homage is that he delivered the needy from the
mighty, and ``this needy and poor man is the people of men believing in Him.
In this people are also kings adoring Him,'' who do not disdain to be poor,
humbly confessing their sins (71.14). So Augustine gives the kings two
places: they belong to the poor people who confess their sins, and their
authority has an office in the Church's spread. On the second psalm he speaks
to the kings directly: ``now that I am set as King, be ye not sad, kings of the
earth, as if your excellency were taken from you, but rather understand and be
instructed. For it is expedient for you, that ye should be under Him''
(2.9).

In his \work{Letter} 185, written to justify the imperial laws against the
Donatists, Augustine himself connects the Offertory's psalm with the
Gradual's verse. In the apostles' time no emperor yet served Christ by laws, because
the raging of the kings in Ps 2:1--2 was being fulfilled and not yet the
command of verses 10--11. ``How then are kings to serve the Lord with fear,
except by preventing and chastising with religious severity all those acts
which are done in opposition to the commandments of the Lord? For a man serves
God in one way in that he is man, in another way in that he is also king''; and
``so soon as the fulfillment began of what is written in a later psalm, `All
kings shall fall down before Him; all nations shall serve Him,' '' no sober man
could tell kings that the Church was no concern of theirs (\work{Letter}
185.19--20). St Thomas reports the same distinction on Ps 2:10--11,
\latin{secundum Augustinum \dots\ rex servit Deo inquantum homo, in se juste
vivendo, sed inquantum rex, leges ferendo contra ea quae sunt contra Dei
justitiam}, and sees in the psalm a figure of the Church's history: at first the
kings of the earth made laws against Christ and Christians, and afterwards laws
for Christ (\work{Postilla super Psalmos} 2, n.~9).

St Robert Bellarmine reads the kings' part as conversion. On the second psalm
the prophet ``exhorts the kings of this world \dots\ to freely subject
themselves to him, the true and supreme king of all kings,'' correcting first
their judgment and then their affections, and bids them serve the Lord as a son
serves his father (\work{Explanatio}, Ps~2). On the Gradual's verse he
distinguishes three fulfilments. The adoration of all kings is fulfilled ``to a
certain extent'' in that ``a great many princes and nations will be converted
to the service and worship of Christ''; it is fulfilled ``to the letter'' at the
day of judgment, when every knee shall bend; and it is fulfilled ``always'' if
it is referred to ``the actual power that Christ has over all princes and all
nations, so that, with or without their knowledge, with or without their
consent, he may deal as he pleases with them'' (Mt 28:18; Apoc 1:5; 19:16). They
serve him, he adds, because ``through him they will be delivered from the power
of the devil \dots\ and will be introduced to his own most peaceful kingdom''
(\work{Explanatio}, Ps~71).

Theodoret divides the same verse in time: some kings adore willingly in the
present life, all after the resurrection, since ``now we see not as yet all
things subject to him'' (Heb 2:8), but then every knee shall bow (Phil 2:10);
and by his laws Christ restrains the powerful and makes them merciful to the
poor (\work{Interpretatio in Psalmos} 71). Cassiodorus reads ``all kings'' as
all nations, lest the verse be thought to speak of kings only, and sharpens the
paradox: \latin{reges adorabunt, qui putantur adorandi; gentes autem servire,
quae terrenos dominos habere noscuntur}, the kings will adore who are thought
to be adored, and the nations serve who are known to have earthly lords
(\work{Expositio Psalmorum} 71). Blessed Ildefonso Schuster's note on
the Gradual gives the state of things between the two fulfilments: the kingdom
is ``complete and absolute de jure, but \dots\ opposed and incomplete de facto''
(p.~202). On the Gospel he writes that the kingdom is not \latin{hinc} ``but it
comprehends the world,'' and that ``Jesus does not come to dethrone the
sovereigns of the earth,'' but gives to rulers and subjects the laws of truth
and justice (p.~203).

\subsection{The kingdom's claim on rulers and Caesar's throne}
\label{sec:hence-difficulty}

The two misreadings set out at the start can each be drawn from Augustine.
His ``I interfere not with your government'' could be pressed to mean that the
kingdom has no public consequence at all, and his \work{Letter} 185 could be
pressed the other way, to make the Gospel a programme of rule. Both statements
are Augustine's, so he cannot have meant either in a sense that excludes the
other.
Chrysostom's reading of the denial lets the two stand together. The denial
before Pilate is a denial of a worldly origin and of worldly means, not of
authority over the world: his empire is from above, and does not deprive the
world of his providence. The subjection the Mass asks for is real, and these
authors give it three moments. It is rendered freely, in the conversion of kings and nations who
come into the poor people that confesses its sins (Augustine, \work{Enarrationes
in Psalmos} 71.14; Bellarmine; Theodoret's ``willingly''). It is rendered by
rulers as rulers, in the exercise of their own office (Augustine, \work{Letter}
185.19; Aquinas). And it is completed at the judgment, when every knee bends
(Bellarmine, Theodoret, Schuster). Beneath all three lies a dominion that
depends on no one's consent, Bellarmine's ``actual power'' and Cyril's King
``whether man so will, or no.'' None of these makes Christ a competitor for
Caesar's throne. \work{Quas primas} says both: Christ during his earthly life refrained from exercising such authority and does not
now interfere with those who possess earthly goods, and yet individuals,
families and states are under his rule, and rulers who know they rule by his
mandate owe him public honour and obedience (nn.~17--19, 32).

Augustine and Chrysostom also differ on a narrower point. For Augustine the
kingdom of Jn 18:36 is the people ruled; for Chrysostom it is the royal power;
Aquinas keeps both. The two divide again on the next verse. Augustine reads
``For this was I born'' of the birth from the Virgin, Chrysostom of the eternal
generation (see the Gospel, above). This interpretation rests on the origin of
the kingdom in verse~36 and on the hearers of the truth in verse~37, and holds
with either birth.

\subsection{The kingdom and the nations, from the Introit to the Postcommunion}
\label{sec:hence-elements}

The Introit acclaims the slain Lamb before the throne, and its verse asks that
God give the King his judgment and his justice; Augustine reads the psalm of a
king whose office is judgment and justice to the poor. Bellarmine, reading the
verse literally of David's prayer for Solomon, says that ``kings were created
for the people, not the people for kings,'' and applied to the kings of the
Gradual, who serve Christ in their own office, that rule means that their
office too is for the people. The Collect says the nations are scattered by
the wound of sin, \latin{peccati vulnere disgregatae}, ``a disruptive force''
in Schuster's note, and asks that they be subjected to a rule it calls
\latin{suavissimum}, most gentle. With Augustine's sense of \latin{regnum},
the people ruled, the petition asks for the nations' gathering into the
Church; with Chrysostom's, the royal power, it asks for their subjection to a
power that no earthly kingdom rivals.

In the Epistle one enters the kingdom by being \latin{translatus}, carried
over out of the power of darkness. Augustine uses the verse to explain the
Gospel, and Aquinas glosses it \latin{ut essemus regnum Dei} with the same
Gospel verse beside it. The Gradual's ``sea to sea'' is, for Augustine, the
Church through the world, and its river is the Jordan, from which, in
Augustine, Cassiodorus and Theodoret, the Gospel and baptism went out to all
nations; its kings are those whose authority leads men into the Church
(Augustine) and who serve Christ as kings (Augustine), whose homage
is given partly now and wholly at the judgment (Bellarmine, Theodoret). The Alleluia sings the power whose origin Aquinas,
with Chrysostom, places above the world, and in Schuster's double chant it
adds to the Gradual's universality of place and subjects the universality of
the ages.

In the Gospel the bound prisoner says his kingdom is not of this world,
\latin{de hoc mundo}, and not from here, \latin{hinc}; he says he was born to
give testimony to the truth, and that every one who is of the truth hears his
voice. Those hearers are the people Augustine calls the kingdom. The Offertory
promises the nations as inheritance, in a psalm whose next verses bid the kings
understand and serve; Bellarmine reads there the kings' free subjection,
Augustine and Aquinas their service as kings. The Secret asks the Son
immolated in the sacrifice to give all nations the gifts of unity and peace.
Set beside the Collect, whose nations are scattered by sin, the unity asked
for is their gathering, and the Secret asks for it as a gift. The Preface names
the kingdom first a kingdom of truth, \latin{regnum veritatis}, and in the
Gospel the King's subjects are those who are of the truth. The Communion's blessing in peace is,
for Augustine, not peace in this world, which the Lord ``hath not promised,''
but peace in himself (\work{Enarrationes in Psalmos} 28.11); Theodoret reads it
of the peace given to the churches after the war the martyrs won.

In the Postcommunion the faithful fight under his standards, \latin{sub
Christi Regis vexillis militare}, after a Gospel in which the King's servants do not fight.
Chrysostom's point on Jn 18:36 is that earthly kingship depends on servants who
fight and Christ's does not; Hilary's, on the second psalm, is that the kingdom
of God is where sin is overcome and death destroyed (\work{Tractatus super
Psalmos} 2.42). Neither speaks of the Postcommunion. Put together, the two
points mean that the faithful fight not to defend his kingdom, which needs no servants who fight,
but against sin, sharing his victory over it. The standard itself, in
Schuster's note on the prayer, is the Cross (p.~204).

\begin{fourSenses}
\item[Literal.] Before the Roman governor Jesus denies that his kingdom has a
worldly origin or worldly defenders, and declares that he was born to bear
witness to the truth. Royal psalms foresee a king whose dominion reaches the
ends of the earth, whom all kings adore and all nations serve, and they warn
the kings of the earth to serve the Lord.
\item[Allegorical.] The kingdom is the Church, here and not from here, gathered
from all nations into one by the truth and spread from the Jordan through the
world (Augustine). Kings enter it as members of the poor people who confess
their sins, and their authority leads others in (Augustine, \work{Enarrationes
in Psalmos} 71.13--14).
\item[Moral.] To be his subject is to be of the truth and to hear his voice.
Every station owes him service in its own office: every man by living
faithfully, the ruler also by ruling justly (Augustine, \work{Letter} 185.19;
\work{Quas primas} nn.~18--19). The warfare of his subjects is against sin,
under the standard of the Cross.
\item[Anagogical.] At the harvest the angels will gather out of his kingdom all
that offends (Mt 13:41, in Augustine's reading). What kings and nations now
give willingly, all will give after the resurrection, when every knee bends
(Theodoret, Bellarmine). Then the kingdom now mixed with tares is made
manifest, and the peace he gives in himself is no longer kept amid the storms
of this world.
\end{fourSenses}
